\chapter*{Lời mở đầu}
\addcontentsline{toc}{chapter}{Lời mở đầu}

Báo cáo cuối kỳ này tập trung vào bài báo \textit{Sample and Computation Redistribution for Efficient Face Detection} (SCRFD) của Guo và cộng sự, công bố tại ICLR 2022. Mục tiêu của báo cáo không chỉ là tóm tắt nội dung paper, mà là trình bày bài báo theo cấu trúc của một báo cáo khoa học hoàn chỉnh: nêu động lực nghiên cứu, phân tích phương pháp, đối chiếu paper với mã nguồn chính thức, tái hiện thực nghiệm và thảo luận ý nghĩa của kết quả thu được.

Từ báo cáo giữa kỳ, nhóm chỉ kế thừa phần liên quan trực tiếp đến Chương II về SCRFD: kiến trúc Backbone--Neck--Head, hai cơ chế \textit{Computation Redistribution} và \textit{Sample Redistribution}, cùng các nhận định về hiệu quả trên WIDER FACE. Phạm vi cuối kỳ được thu hẹp về một bài báo khoa học cụ thể và phần tái hiện thực nghiệm của bài báo đó.

Phần tái hiện chính sử dụng mã nguồn \textit{SCRFD} trong cây \texttt{insightface/detection/scrfd}, cùng các script thí nghiệm ở nhánh \texttt{experiments} để tổ chức lại các run baseline, paper-SR12 và phần mở rộng \textit{ASR+JSAR}. Nhánh SPOS-NAS proxy được giữ lại như một phân tích phụ để minh họa thêm trực giác về tái phân bổ tính toán, nhưng không thay thế cho phần tái hiện chính bằng official code.

Bố cục báo cáo gồm bảy chương. Chương I giới thiệu bài báo SCRFD và bài toán nghiên cứu. Chương II trình bày phương pháp trong paper. Chương III mô tả quá trình tái hiện mã nguồn chính thức. Chương IV đối chiếu các ý tưởng trong paper với module code tương ứng. Chương V trình bày kết quả thực nghiệm và so sánh với kết quả công bố. Chương VI thảo luận các phần mở rộng. Chương VII tổng kết đóng góp và nêu rõ các nội dung đã hoàn thiện so với bản giữa kỳ.

Điều nhóm chúng em thấy giá trị nhất sau quá trình này là: với các hệ detector hiện đại, hiểu bài báo thôi chưa đủ; cần chỉ ra được cấu hình nào, assigner nào, pipeline augmentation nào và script đánh giá nào đã biến ý tưởng trong paper thành hành vi quan sát được trong code. Chính bước nối này làm rõ sự khác nhau giữa ``đọc hiểu bài báo'' và ``thực sự tái hiện được phương pháp''.

\vspace{1.5cm}
\begin{flushright}
\textit{Thành phố Hồ Chí Minh, ngày 19 tháng 4 năm 2026}\\[0.5cm]
\textbf{Nhóm sinh viên thực hiện}
\end{flushright}
